\section{Bounded ordered sets}\label{sec:bounded_ordered_sets}

\paragraph{Extremal points}

\begin{definition}\label{def:extremal_points}
  We introduce some notions for extremal elements of a \hyperref[def:partially_ordered_set]{partially ordered set} \( (P, \leq) \). Bounds, introduced first, do not rely on antisymmetry and thus generalize to \hyperref[def:preordered_set]{preordered sets}, but the other notions can be ill-behaved (e.g. \hyperref[def:gcd]{greatest common divisors} are not unique in general).

  \begin{thmenum}
    \thmitem{def:extremal_points/bounds}
    \begin{paracol}{2}
      \begin{leftcolumn}\mcite[def. 1.6.6]{Harzheim2005OrderedSets}
        An \term[bg=горна граница (\cite[18]{Тагамлицки1971ДиференциалноСмятане}) / мажоранта (\cite[10]{Проданов1982ФункционаленАнализЧаст1}), ru=верхняя грань / мажоранта (\cite[def. 3.8]{Гуров2013ТеорияРешёток})]{upper bound} for the set \( A \) is an element \( m \) of the ambient ordered set such that \( x \leq m \) for every \( x \) in \( A \).

        We say that the upper bound is \term{strict} if the inequality is strict, i.e. if \( m \) does not itself belong to \( A \).

        Nonstrict upper bounds generalize to preordered sets.

        If \( A \) has at least one upper bound, we say that \( A \) is \term{bounded from above}.
      \end{leftcolumn}

      \begin{rightcolumn}
        Dually, a \term[bg=долна граница (\cite[18]{Тагамлицки1971ДиференциалноСмятане}) / миноранта (\cite[10]{Проданов1982ФункционаленАнализЧаст1}), ru=нижняя грань / миноранта (\cite[def. 3.8]{Гуров2013ТеорияРешёток})]{lower bound} for \( A \) is an element \( m \) such that \( m \leq x \) for every \( x \) in \( A \).

        If \( A \) has a lower bound, we say that \( A \) is \term{bounded from below}.

        If \( A \) is both bounded from above and from below, we simply say that it is \term[bg=ограничено (множество) (\cite[19]{Тагамлицки1971ДиференциалноСмятане}), ru=ограниченное (множество) (\cite[39]{Зорич2019АнализЧасть1})]{bounded}.

        Boundedness also generalizes to preordered sets.
      \end{rightcolumn}
    \end{paracol}

    \thmitem{def:extremal_points/maximal_and_minimal_element}
    \begin{paracol}{2}
      \begin{leftcolumn}\mcite[def. III.1.1.3]{Harzheim2005OrderedSets}
        A \term[ru=максимальный (елемент) (\cite[def. 3.6]{Гуров2013ТеорияРешёток})]{maximal element} for the set \( A \subseteq P \) is a member \( m \) of \( A \) such that \( x \geq m \) implies \( x = m \) for every \( x \) in \( A \).

        Equivalently, we can require \( x \leq m \) whenever \( x \) and \( m \) are comparable (and \( x \) is in \( A \)).
      \end{leftcolumn}

      \begin{rightcolumn}
        Dually, \( m \) is \term[ru=минимальный (елемент) (\cite[def. 3.6]{Гуров2013ТеорияРешёток})]{minimal} if \( x \leq m \) implies \( x = m \) for every \( x \) in \( A \).

        Equivalently, we can require \( x \geq m \) whenever \( x \) and \( m \) are comparable (and \( x \) is in \( A \)).
      \end{rightcolumn}
    \end{paracol}

    \thmitem{def:extremal_points/maximum_and_minimum}
    \begin{paracol}{2}
      \begin{leftcolumn}\mcite[def. 1.6.8]{Harzheim2005OrderedSets}
        If two upper bounds \( m \) and \( m' \) of \( A \) both belong to \( A \), they must be equal because \( m \leq m' \) and \( m' \leq m \) implies \( m = m' \). We call this unique element the \term[ru=наибольший (елемент) (\cite[def. 3.6]{Гуров2013ТеорияРешёток})]{greatest element} or \term{maximum} and denote it by \( \max A \).
      \end{leftcolumn}

      \begin{rightcolumn}
        Dually, if a lower bound of \( A \) belongs to \( A \), we call it the \term[ru=наименьший (елемент) (\cite[def. 3.6]{Гуров2013ТеорияРешёток})]{least element} or \term{minimum} of \( A \) and denote it by \( \min A \).
      \end{rightcolumn}
    \end{paracol}

    \thmitem{def:extremal_points/supremum_and_infimum}
    \begin{paracol}{2}
      \begin{leftcolumn}\mcite[def. 1.6.8]{Harzheim2005OrderedSets}
        If the upper bounds of \( A \) have a minimum, we call it the \term[bg=точна горна граница (\cite[19]{Тагамлицки1971ДиференциалноСмятане}), ru=точная верхняя граница (\cite[\S 5.12]{Ляпин1960Полугруппы})]{least upper bound} or \term[bg=супремум (\cite[10]{Проданов1982ФункционаленАнализЧаст1}), ru=супремум (\cite[\S 5.12]{Ляпин1960Полугруппы})]{supremum} of \( A \) and denote it by \( \sup A \).
      \end{leftcolumn}

      \begin{rightcolumn}
        Dually, if the lower bounds of \( A \) have a maximum, we call it the \term[bg=точна долна граница (\cite[19]{Тагамлицки1971ДиференциалноСмятане}), ru=точная нижняя граница (\cite[\S 5.12]{Ляпин1960Полугруппы})]{greatest lower bound} or \term[bg=инфимум (\cite[10]{Проданов1982ФункционаленАнализЧаст1}), ru=супремум (\cite[\S 5.12]{Ляпин1960Полугруппы})]{infimum} of \( A \) and denote it by \( \inf A \).
      \end{rightcolumn}
    \end{paracol}

    \thmitem{def:extremal_points/top_and_bottom}
    \begin{paracol}{2}
      \begin{leftcolumn}\mcite[15]{DaveyPriestley2002LatticeTheory}
        If the ambient space itself has a maximum, we call it the \term{top element} and denote it via \( \top \).
      \end{leftcolumn}

      \begin{rightcolumn}
        Dually, if the ambient space has a minimum, we call it the \term{bottom element} and denote it by \( \bot \).
      \end{rightcolumn}
    \end{paracol}
  \end{thmenum}
\end{definition}

\begin{proposition}\label{thm:def:extremal_points}
  Fix a \hyperref[def:partially_ordered_set]{partially ordered set} \( P \) and some nonempty subset \( A \) of \( P \). The extremal point notions from \cref{def:extremal_points} have the following basic properties:
  \begin{thmenum}
    \thmitem{thm:def:extremal_points/empty_exact_bounds} If \( P \) has a bottom element, it is the supremum of the empty set.

    Dually, if \( P \) has a top element, it is the infimum of the empty set.

    \thmitem{thm:def:extremal_points/greatest_is_maximal} The greatest element of \( A \), if it exists, is also a maximal element of \( A \).

    Dually, the least element of \( A \) is a minimal element of \( A \).

    The converse holds in \hyperref[def:totally_ordered_set]{totally ordered sets} --- see \cref{thm:def:totally_ordered_set/maximal_iff_greatest}.

    \thmitem{thm:def:extremal_points/greatest_is_supremum} The greatest element of \( A \), if it exists, is the supremum of \( A \).

    Dually, the least element of \( A \) is the infimum of \( A \).

    \thmitem{thm:def:extremal_points/finite_set_has_maximal_element} If \( A \) is \emph{finite}, it has both a maximal element and a minimal element.
  \end{thmenum}
\end{proposition}
\begin{proof}
  \SubProofOf{thm:def:extremal_points/empty_exact_bounds} A supremum is a minimum among all upper bounds. All elements of \( P \) are vacuously upper bounds of \( \varnothing \) since there is nothing to compare them to, so a minimum of \( P \) is a supremum of \( \varnothing \).

  \SubProofOf{thm:def:extremal_points/greatest_is_maximal} If \( m \) is the greatest element of \( A \), then, by definition, we have \( m \geq x \) for every \( x \) in \( A \).

  If \( x \geq m \) for some \( x \) in \( A \), from antisymmetry it follows that \( x = m \). So \( m \) is a maximal element of \( A \).

  \SubProofOf{thm:def:extremal_points/greatest_is_supremum} Similarly, let \( m \) be the greatest element of \( A \).

  Then \( m \) is by definition an upper bound of \( A \), but, since \( m \) belongs to \( A \), we have \( m \leq u \) for every upper bound \( u \) of \( A \). So \( m \) is the least upper bound of \( A \), i.e. the supremum.

  \SubProofOf{thm:def:extremal_points/finite_set_has_maximal_element} Suppose that \( A \) is finite with elements
  \begin{equation*}
    a_1, a_2, \ldots, a_n.
  \end{equation*}

  We will use induction on \( n \) to show that \( A \) has a maximal element.
  \begin{itemize}
    \item The base case \( n = 1 \) is trivial --- \( A \) has one element, which can only be maximal.

    \item Suppose that every set of size \( n - 1 \) has a maximal element. Let \( a_m \) be a maximal element of \( A \setminus \set{ a_n } \) . Then we have the following possibilities:
    \begin{itemize}
      \item If \( a_m \) is not comparable to \( a_n \), then \( a_m \) is also maximal for \( A \).

      \item If \( a_m < a_n \), then \( a_n \) is maximal for \( A \) because \( a_k \geq a_n \) can only hold when \( k = n \).

      \item If \( a_m \geq a_n \), the \( m \) is maximal for \( A \) because \( a_k \geq a_m \) implies that \( a_k = a_m \) in case \( k < n \) and that \( a_n \geq a_m \) in case \( k = n \).
    \end{itemize}

    In all cases, we have shown that \( A \) has at least one maximal element.
  \end{itemize}

  We can similarly show that \( A \) has a minimal element.
\end{proof}

\begin{example}\label{ex:thm:def:extremal_points}
  We list examples of for the extremal point notions from \cref{def:extremal_points}:
  \begin{thmenum}
    \thmitem{ex:thm:def:extremal_points/empty} Fix some arbitrary ordered set \( P \) and consider the empty subset \( \varnothing \).

    Clearly every member of \( P \) is both an upper and a lower bound of \( \varnothing \). If \( P \) has a bottom element, it is a supremum of \( \varnothing \) because it is a minimum among all upper bounds of \( \varnothing \), that is, among all elements of \( P \). Similarly, if \( P \) has a top element, it is an infimum of \( \varnothing \).

    Since \( \varnothing \) is empty, it cannot have a greatest element, nor a maximal element.

    \thmitem{ex:thm:def:extremal_points/one} Consider the one-element set \( \set{ a } \). The only possible reflexive relation entails \( a \leq a \), so \( a \) is the top and bottom of \( \set{ a } \).

    \thmitem{ex:thm:def:extremal_points/two_isolated} Consider the two-element set \( \set{ a, b } \) with the \hyperref[def:binary_relation/diagonal]{diagonal relation} (every element is related only to itself).

    Then \( a \) is the only upper bound of \( \set{ a } \), hence also a maximal element, greatest element and supremum. It is similarly a lower bound, minimal element, minimum and infimum.

    But for \( \set{ a, b } \), it is neither an upper nor lower bound.

    On the other hand, both \( a \) and \( b \) are maximal elements of \( \set{ a, b } \), because there exist no elements greater than them. Similarly, they are both minimal.

    \thmitem{ex:thm:def:extremal_points/two_connected} Consider again the two-element set \( \set{ a, b } \), but this time with \( a \leq b \). Here \( a \) is a bottom element and \( b \) is a top element.

    All such ordered sets are necessarily isomorphic.

    \thmitem{ex:thm:def:extremal_points/two_equivalent} Consider yet again the two-element set \( \set{ a, b } \), but this time with both \( a \leq b \) and \( b \leq a \). Then \( {\leq} \) is not antisymmetric, hence it is a preorder.

    Both are top elements and bottom elements because of the lack of uniqueness required by the definitions of maximum and minimum.

    \thmitem{ex:thm:def:extremal_points/three_incomparable} Consider the three-element set \( \set{ a, b, c } \) such that \( a \leq c \) and \( b \leq c \).

    Clearly \( c \) is a top element. But there is no bottom. Both \( a \) and \( b \) are minimal elements of \( \set{ a, b, c } \), but neither of them is a minimum because they are not comparable.

    \thmitem{ex:thm:def:extremal_points/unique_maximal_element_not_greatest}\mcite{MathSE:unique_maximal_element_that_is_not_greatest} Consider the set \( \BbbZ \) of integers with the standard ordering from \cref{def:integer_ordering}.

    Adjoin some symbol \( u \) and define a relation to \( \BbbZ \cup \set{ u } \) extending the ordering on \( \BbbZ \) with \( u \leq u \).

    It is the unique maximal element of \( \BbbZ \cup \set{ u } \) because no integer is greater than \( u \). But, since it is incomparable with integers, \( u \) is not a greatest element.

    \thmitem{ex:thm:def:extremal_points/nonunique_maxima} Consider again the integers \( \BbbZ \), but this time adjoin two incomparable elements --- \( \infty \) and \( \tieinfty \) --- and let both of them be greater than any integer.

    Then both \( \infty \) and \( \tieinfty \) are upper bounds of \( \BbbZ \) in \( \BbbZ \cup \set{ \infty, \tieinfty } \). Furthermore, both of them are minimal in the set \( \set{ \infty, \tieinfty } \) of upper bounds, but, since they are not comparable, neither is a least upper bound.
  \end{thmenum}
\end{example}

\paragraph{Intervals}

\begin{definition}\label{def:order_interval}\mcite[def. 1.8.9]{Harzheim2005OrderedSets}
  Fix a \hyperref[def:partially_ordered_set]{partially ordered set} \( (P, \leq) \). For any two elements \( a \leq b \) we define the following related (ordered) subsets:

  \begin{thmenum}
    \thmitem{def:order_interval/unbounded} First, we define the \term{open initial segment} and \term{open final segment} and as
    \begin{equation*}
      \begin{aligned}
        P_{>a} \coloneqq \set{ x \in P \given b > a },
        \\
        P_{<b} \coloneqq \set{ x \in P \given x < b }.
      \end{aligned}
    \end{equation*}

    Similarly, we define the corresponding \term{closed segments} as
    \begin{equation*}
      \begin{aligned}
        P_{\geq a} \coloneqq \set{ x \in P \given x \geq a },
        \\
        P_{\leq b} \coloneqq \set{ x \in P \given x \leq b }.
      \end{aligned}
    \end{equation*}

    \thmitem{def:order_interval/closed} Next, we define the \term[bg=затворен интервал (от реални числа) (\cite[39]{Тагамлицки1971ДиференциалноСмятане}), ru=отрезок (\cite[82]{Проскуряков1951ТеоретическиеОсновыАрифметики}); отрезок/сегмент (\cite[def. 6]{Александров1977ОбщаяТопология})]{closed interval} with endpoints \( a \) and \( b \) as
    \begin{equation*}
      [a, b] \coloneqq \set{ x \in P \given a \leq x \leq b } = P_{\geq a} \cap P_{\leq b}.
    \end{equation*}

    \thmitem{def:order_interval/open} Next, we define the \term[bg=отворен интервал (от реални числа) (\cite[39]{Тагамлицки1971ДиференциалноСмятане}), ru=интервал (\cite[82]{Проскуряков1951ТеоретическиеОсновыАрифметики}) (действительных чисел); интервал/промежуток (\cite[def. 6]{Александров1977ОбщаяТопология})]{open interval} with endpoints \( a \) and \( b \) as
    \begin{equation*}
      (a, b) \coloneqq \set{ x \in P \given a < x < b } = P_{> a} \cap P_{< b}.
    \end{equation*}

    \thmitem{def:order_interval/half_open} Finally, we define the \term[bg=полузатворен интервал (\cite[39]{Тагамлицки1971ДиференциалноСмятане}), ru=полуинтервал (действительных чисел) (\cite[82]{Проскуряков1951ТеоретическиеОсновыАрифметики}), en=half-open interval (of real numbers) (\cite[def. 2.17]{Rudin1976AnalysisPrinciples})]{half-open intervals} are
    \begin{equation*}
      \begin{aligned}
        (a, b] \coloneqq \set{ x \in P \given a < x \leq b } = P_{> a} \cap P_{\leq b},
        \\
        [a, b) \coloneqq \set{ x \in P \given a \leq x < b } = P_{\geq a} \cap P_{< b}.
      \end{aligned}
    \end{equation*}
  \end{thmenum}
\end{definition}
\begin{comments}
  \item We implicitly assume that \( a \leq b \), but this is not strictly necessary --- \( [a, b] \) is an empty set otherwise.

  \item In the context of \hyperref[def:extended_real_numbers]{extended real numbers}, we use \( -\infty \) and \( \infty \) as interval endpoints. If the interval is open at \( -\infty \) and \( \infty \), as in \( (-\infty, b] = \BbbR_{\leq b} \) or \( [a, \infty) = \BbbR_{\geq a} \), we obtain sets containing only real numbers.

  \item Another notation for open intervals is \enquote{\( ]a, b[ \)}. It is used in
  \cite[147]{Bourbaki2004TheoryOfSets},
  \cite[def. 2.3]{EkelandTemam1976ConvexAnalysis} and
  \cite[51]{Зорич2019АнализЧасть1}.

  We prefer \enquote{\( (a, b) \)}, which is in turn used by most of our sources ---
  \cite[9]{Lorentz1966ApproximationOfFunctions},
  \cite[31]{Rudin1976AnalysisPrinciples},
  \cite[10]{Engelking1989GeneralTopology},
  \cite[269]{Jacobson1985BasicAlgebraI},
  \cite[157]{Grätzer2011LatticeTheory},
  \cite[def. 1.8.9]{Harzheim2005OrderedSets},
  \cite[\S 3.3.14]{Schechter1997AnalysisHandbook}
  \cite[11]{Folland1999RealAnalysis}
  \cite[25]{МагарилИльяевТихомиров2002ВыпуклыйАнализ},
  \cite[18]{ИльинСадовничийСендов1985АнализТом1},
  \cite[\S 22.2]{Тыртышников2007ЛинейнаяАлгебра},
  \cite[12]{Александров1977ОбщаяТопология},
  \cite[112]{Винберг2014КурсАлгебры},
  \cite[10]{БелоусовТкачёв2004ДискретнаяМатематика},
  \cite[23]{Гуров2013ТеорияРешёток},
  \cite[7]{Боянов2008ЧислениМетоди},
  \cite[225]{ГеновМиховскиМоллов1991Алгебра} and
  \cite[39]{Тагамлицки1971ДиференциалноСмятане}.

  \Textcite[36]{Knuth1997ArtVol1} uses \enquote{\( (a..b) \)}.
\end{comments}

\paragraph{Zorn's lemma}

\begin{lemma}[Zorn's lemma]\label{thm:zorns_lemma}
  If every \hyperref[def:partial_order_connectedness]{chain} in a nonempty \hyperref[def:partially_ordered_set]{partially ordered set} has an \hyperref[def:extremal_points/bounds]{upper bound}, then the entire set has a \hyperref[def:extremal_points/maximal_and_minimal_element]{maximal element}.
\end{lemma}
\begin{comments}
  \item Within \hyperref[def:zfc]{\logic{ZF}}, this theorem is equivalent to the \hyperref[def:zfc/choice]{axiom of choice} --- see \cref{thm:axiom_of_choice_equivalences/zorns_lemma}.

  \item Zorn's lemma is sometimes stated and used only in a \hyperref[thm:boolean_algebra_of_subsets]{lattice of sets} --- for example in
  \cite[117]{Grätzer2011LatticeTheory},
  \cite[thm. 6.4.34]{Hinman2005Logic},
  \cite[151]{Enderton1977SetTheory} and
  \cite[16]{КанторовичАкилов1984ФункциональныйАнализ}.

  A statement similar to this one is used by
  \cite[50]{Harzheim2005OrderedSets},
  \cite[33]{Kelley1975GeneralTopology},
  \cite[8]{Engelking1989GeneralTopology},
  \cite[lemma V.3.1]{Aluffi2009Algebra},
  \cite[880]{Lang2002Algebra},
  \cite[314]{Rotman2015AdvancedModernAlgebraPart1},
  \cite[117]{Гуров2013ТеорияРешёток},
  \cite[31]{Мальцев1970АлгебраическиеСистемы} and
  \cite[13]{Проданов1982ФункционаленАнализЧаст1}.

  \Textcite[2]{Jacobson1989BasicAlgebraII} states both.
\end{comments}
\begin{proof}
  \ImplicationSubProof[def:zfc/choice]{the axiom of choice}[thm:zorns_lemma]{Zorn's lemma} Let \( (P, \leq) \) be a partially ordered set in which every chain has an upper bound. Aiming at a contradiction, suppose that \( P \) has no maximal elements.

  Let \( \mscrC \) be the set of all chains in \( P \). Define the set-valued map \( F: \mscrC \multto P \) that assigns to each set in \( P \) the family of all its strict upper bounds.

  Since every chain \( C \) has an upper bound, say \( u \), and since \( P \) has no maximal element, there exists some element strictly larger than \( u \). Hence, \( F \) is a total set-valued map. By \fullref{thm:existence_of_single_valued_selections}, there exists a single-valued selection \( f: \pow(P) \to P \) of \( F \).

  By \fullref{thm:hartogs_lemma}, there exists a smallest \hyperref[def:ordinal]{ordinal} \( \alpha \) such that no function from \( \alpha \) to \( P \) is injective. Using \fullref{thm:bounded_transfinite_recursion}, we can define
  \begin{equation*}
    \begin{aligned}
      &g: \alpha \to P, \\
      &g(\beta) \coloneqq f(\set{ g(\gamma) \given \gamma < \beta }) = f(g[\beta]). \\
    \end{aligned}
  \end{equation*}

  If \( \beta_1 < \beta_2 \), then
  \begin{equation*}
    \set{ g(\gamma) \given \gamma < \beta_1 }
    \subsetneq
    \set{ g(\gamma) \given \gamma < \beta_2 },
  \end{equation*}
  and thus
  \begin{equation*}
    g(\beta_1) < g(\beta_2).
  \end{equation*}

  Then the images under \( g \) of different elements of \( \alpha \) are different, i.e. \( g \) is injective. But this contradicts our choice of \( \alpha \).

  The obtained contradiction shows that \( P \) has a maximal element.

  \ImplicationSubProof[thm:zorns_lemma]{Zorn's lemma}[def:zfc/choice]{axiom of choice} Let \( \mscrA \) be a family of nonempty sets. Let \( \mscrF \) be the set of all \hyperref[def:set_valued_map/partial]{partial single-valued functions} from \( \mscrA \) to \( \bigcup \mscrA \) with the subset ordering. That is, \( f \leq g \) if \( \dom(f) \subseteq \dom(g) \) for \( f, g \in \mscrF \).

  Clearly every chain has a greatest element - a total single-valued function. Then \( \mscrF \) itself has a maximal element by Zorn's lemma. This maximal element is necessarily a total function because otherwise it would not be maximal.

  Then this is the desired choice function for the family \( \mscrA \).
\end{proof}

\paragraph{Chains}

\begin{definition}\label{def:partial_order_connectedness}\mcite[4]{Grätzer2011LatticeTheory}
  Fix a partially ordered set \( (P, \leq) \) and a subset \( C \) of \( P \). We will study the extremes regarding the \hyperref[def:binary_relation/connected]{connectedness} of \( C \).

  \begin{thmenum}
    \thmitem{def:partial_order_connectedness/chain} We say that \( C \) is a \term[bg=верига (\cite[10]{Проданов1982ФункционаленАнализЧаст1}), ru=цепь (\cite[def. 3.4]{Гуров2013ТеорияРешёток})]{chain} if every pair of elements of \( C \) is comparable.

    \thmitem{def:partial_order_connectedness/length}\mimprovised We define the \term[ru=длина (\cite[83]{Гуров2013ТеорияРешёток})]{length} of a \emph{nonempty} chain \( C \) as \( \len(C) = \card(C) - 1 \), i.e. the unique cardinal satisfying \( \len(C) + 1 = \card(A) \).

    When \( C \) is infinite, we have \( \len(C) = \card(C) \).

    \thmitem{def:partial_order_connectedness/height}\mimprovised We define the \term[ru=высота (\cite[83]{Гуров2013ТеорияРешёток})]{height} of a \emph{nonempty} ordered set as the supremum of the lengths of its nonempty chains.

    \thmitem{def:partial_order_connectedness/antichain} Dually to chains, we say that \( C \) is an \term[ru=антицепь (\cite[78]{Гуров2013ТеорияРешёток})]{antichain} if no two elements of \( C \) are comparable.

    \thmitem{def:partial_order_connectedness/width} We define the \term[ru=ширина (\cite[83]{Гуров2013ТеорияРешёток})]{width} of a \emph{finite} ordered set as the supremum of the cardinalities of its antichains.
  \end{thmenum}
\end{definition}
\begin{comments}
  \item We order chains and antichains by inclusion, so phrases like \enquote{maximal chain} have a precise meaning.

  \item If we interpret \( (P, \leq) \) as a \hyperref[def:directed_graph]{directed graph} in accordance with \cref{rem:logical_theory_of_graphs/directed_graphs}, finite chains become \hyperref[def:graph_walk/directed]{directed walks}, and in such case the length of a chain, per Gr\"atzer's definition, coincides with the length of the corresponding graph walk, per \cref{def:graph_walk/length}.

  \item There is a certain asymmetry in the definitions of width and height. We discuss it in \cref{rem:partial_order_height_and_width} and \cref{ex:def:partial_order_connectedness}.
\end{comments}

\begin{remark}\label{rem:partial_order_height_and_width}
  Chains, as defined in \cref{def:partial_order_connectedness}, are similarly defined in
  \cite[4]{Grätzer2011LatticeTheory},
  \cite[2]{Birkhoff1967LatticeTheory},
  \cite[def. 1.6.1]{Harzheim2005OrderedSets} and
  \cite[\S 1.3]{DaveyPriestley2002LatticeTheory}. Antichains are also defined by all aforementioned authors except Birkhoff.

  Chains and antichains are additionally assumed to be nonempty in \cite[def. 3.4]{Гуров2013ТеорияРешёток}.

  The length of a \emph{finite} \emph{nonempty} chain \( C \) is defined as \( \card(C) - 1 \) in all aforementioned books except in \cite[def. 1.6.1]{Harzheim2005OrderedSets}, where it is defined as \( \card(C) \). No equivalent concept is defined for antichains.

  Regarding the \enquote{height} (sometimes called the \enquote{length}) and \enquote{width} of a partially ordered set \( (P, \leq) \), all aforementioned authors give slightly different definitions, which coincide when \( P \) is nonempty and finite but are ambiguous more generally. Note that \textcite{Dilworth1950DecompositionTheorem} and \textcite{Mirsky1971DualOfDilworthsTheorem}, whose theorems the width and height are often associated with, do not use these terms. The apparent asymmetry of the width and height is shown to be somewhat superficial in \cref{ex:def:partial_order_connectedness}.

  \begin{itemize}
    \item \Textcite[def. 2.5.1]{Harzheim2005OrderedSets} defines the \enquote{width} of \( P \) as the \enquote{supremum of the cardinals of antichains in \( P \)}.

    We take this definition, and use it as a prototype for the \enquote{height} of \( P \). To contrast, \textcite[def. 2.7.1]{Harzheim2005OrderedSets} defines the \enquote{height} of a \emph{nonempty} ordered set \( P \) as \enquote{the maximum of all numbers \( h(a) + 1 \), \( a \in P \)}, where the \enquote{height} \( h(a) \) of an element \( a \) is \enquote{the greatest nonnegative integer \( h \), so that there exists a chain \( \set{ a_0, \ldots, a_h = a } \), where \( a_0 < \ldots < a_h \)}.

    Harzheim's definition of height is sound, however less general than the definition of width he provides.

    \item \Textcite[4]{Grätzer2011LatticeTheory} defines the \enquote{length} of \( P \) as the natural number \( n \) \enquote{if there is a chain in \( P \) of length \( n \) and all chains in \( P \) are of length \( \leq n \)}.

    Gr\"atzer does not mention what happens if \( P \) is empty or if it contains infinite chains.

    He defines the \enquote{width} of \( P \) dually.

    \item \Textcite[2]{Birkhoff1967LatticeTheory} defines the \enquote{length} of \( P \) as \enquote{the least upper bound of the lengths of the chains in \( P \)}.

    Birkhoff also does not mention what happens if \( P \) is empty, nor does he define the length of an infinite chain.

    He does not discuss antichains.

    \item \Textcite[2.37(ii)]{DaveyPriestley2002LatticeTheory} define the \enquote{length} of \( P \) as \( n \) if \enquote{the length of the longest chain in \( P \) is \( n \)}.

    The length is again ambiguous if \( P \) is empty, and if \( P \) is infinite a longest chain may not exist.

    Nearly the same definition is given in \cite[83]{Гуров2013ТеорияРешёток}, but is called the \enquote{\textru{высота}} (\enquote{height}) of \( P \).

    They define the \enquote{width} (resp. \enquote{\textru{ширина}}) of \( P \) dually.
  \end{itemize}
\end{remark}

\begin{example}\label{ex:def:partial_order_connectedness}
  We list some examples of \hyperref[def:partial_order_connectedness]{chains and antichains}:
  \begin{thmenum}
    \thmitem{ex:def:partial_order_connectedness/natural_numbers} The set \( \BbbN \) of \hyperref[def:natural_numbers]{natural numbers}, as well as any subset of \( \BbbN \), is a chain.

    The length of the initial segment \( \BbbN_{\leq n} = \set{ 0, 1, \ldots, n } \) is \( n \) because it has \( n + 1 \) elements. The length of \( \BbbN \) itself is \( \aleph_0 \).

    Hence, the height of \( \BbbN \) is \( \aleph_0 \), while its width is \( 1 \).

    \thmitem{ex:def:partial_order_connectedness/binary_power_set} Consider the \hyperref[def:basic_set_operations/power_set]{power set} of the two-element set \( \set{ a, b } \) depicted in \cref{fig:ex:def:partial_order_connectedness/binary_power_set}.

    \begin{figure}
      \hfill
      \includegraphics[page=1]{output/ex__def__partial_order_chain__binary_power_set}
      \hfill\hfill
      \caption{A Hasse diagram of the power set of a two-element set.}
      \label{fig:ex:def:partial_order_connectedness/binary_power_set}
    \end{figure}

    Only the sets \( \set{ a } \) and \( \set{ b } \) are incomparable, thus the width of the power set is \( 2 \).

    On the other hand, the maximal chain \( \varnothing \subseteq \set{ a } \subseteq \set{ a, b } \) has three elements, hence is of length two. Thus, the height of the power set is also \( 2 \).

    This shows that the asymmetry in the definition of width and height is more superficial than it seems at first glance.

    \thmitem{ex:def:partial_order_connectedness/ternary_power_set} Consider the \hyperref[def:basic_set_operations/power_set]{power set} of \( \set{ a, b, c } \) depicted in \cref{fig:ex:def:partial_order_connectedness/ternary_power_set}.

    \begin{figure}
      \hfill
      \includegraphics[page=1]{output/ex__def__partial_order_chain__ternary_power_set}
      \hfill\hfill
      \caption{A Hasse diagram of the power set of a three-element set.}
      \label{fig:ex:def:partial_order_connectedness/ternary_power_set}
    \end{figure}

    Like in \cref{ex:def:partial_order_connectedness/binary_power_set}, we can find several maximal antichains:
    \begin{itemize}
      \item The singleton sets \( \set{ a } \), \( \set{ b } \) and \( \set{ c } \).
      \item The binary sets \( \set{ a, b } \), \( \set{ a, c } \) and \( \set{ b, c } \).
    \end{itemize}

    Both have cardinality three, hence the width of the power set is \( 3 \).

    Maximal chains have four elements, so the height of the power set is also \( 3 \).

    \thmitem{ex:def:partial_order_connectedness/quad_power_set} The pattern from \cref{ex:def:partial_order_connectedness/binary_power_set} and \cref{ex:def:partial_order_connectedness/ternary_power_set} breaks when we consider four elements.

    Indeed, for the power set of \( \set{ a, b, c, d } \), the following is an antichain of six elements:
    \begin{equation}\label{eq:ex:def:partial_order_connectedness/quad_power_set/antichain}
      \set{ a, b }, \set{ a, c }, \set{ a, d }, \set{ b, c }, \set{ b, d }, \set{ c, d }.
    \end{equation}
  \end{thmenum}
\end{example}

\begin{proposition}\label{thm:def:partial_order_connectedness}
  \hyperref[thm:def:partial_order_connectedness]{Chains and antichains} have the following basic properties:
  \begin{thmenum}
    \thmitem{thm:def:partial_order_connectedness/image_of_chain} Let \( f: P \to Q \) be an order-preserving map between partially ordered sets. If \( A \) is a chain in \( P \), then its image \( f[A] \) is a chain in \( Q \).

    \thmitem{thm:def:partial_order_connectedness/preimage_of_antichain} Again, let \( f: P \to Q \) be an order-preserving map between partially ordered sets. If the image \( f[A] \) of a subset \( A \) of \( P \) is an antichain, then \( A \) itself is an antichain.

    \thmitem{thm:def:partial_order_connectedness/unbounded} A chain is \hyperref[def:extremal_points/bounds]{unbounded from above} (resp. below) if and only if, for every element, it contains a strictly larger (resp. smaller) element.

    \thmitem{thm:def:partial_order_connectedness/maximal_element_antichain} The set of \hyperref[def:extremal_points/maximal_and_minimal_element]{maximal} (resp. minimal) elements in an ordered set is an \hyperref[def:partial_order_connectedness/antichain]{antichain}.
  \end{thmenum}
\end{proposition}
\begin{proof}
  \SubProofOf{thm:def:partial_order_connectedness/image_of_chain} Fix two elements \( x \) and \( y \) in \( f[A] \). Let \( a \) and \( b \) be elements of \( P \) such that \( f(a) = x \) and \( f(b) = y \). Then \( a \leq_P b \) implies \( f(a) \leq_Q f(b) \) and \( a >_P b \) implies \( f(a) \geq_Q f(b) \). Hence, \( x \) and \( y \) are comparable, and since they were arbitrary, we conclude that \( f[A] \) is a chain.

  \SubProofOf{thm:def:partial_order_connectedness/preimage_of_antichain} Follows from \cref{thm:def:order_function/comparables_reflect_inequalities}.

  \SubProofOf{thm:def:partial_order_connectedness/unbounded} Fix a chain \( C \) in a partially ordered set \( (P, \leq) \).

  Negating \cref{def:extremal_points/bounds}, we obtain
  \begin{displayquote}
    \( C \) is unbounded from above if and only if there does not exist an element \( x \) in \( C \) such that, for every \( y \) in \( C \), we have \( y \leq x \).
  \end{displayquote}

  The consequent of the above statement has a negation in front of its outermost quantifier. \Fullref{alg:fol_formula_dualization} justifies moving the negation inwards to obtain
  \begin{displayquote}
    \( C \) is unbounded from above if and only if, for every \( x \) in \( C \), there exists some \( y \) in \( C \) such that \( y \leq x \) does not hold.
  \end{displayquote}

  Taking into account that every pair of elements in \( C \) is comparable, \( y \leq x \) does not hold if and only if \( y > x \) holds. Then
  \begin{displayquote}
    \( C \) is unbounded from above if and only if, for every \( x \) in \( C \), there exist some \( y \) in \( C \) such that \( y > x \).
  \end{displayquote}

  The case of lower bounds is analogous.

  \SubProofOf{thm:def:partial_order_connectedness/maximal_element_antichain} Let \( M \) be the set of maximal elements of \( (P, \leq) \).

  For any pair \( x \) and \( y \) from \( M \), if \( x \geq y \), by definition of maximal element we have \( x = y \), and similarly if \( y \geq x \).

  Thus, no two distinct incomparable elements exist in \( M \). So it is an antichain.

  If \( M \) consists of the minimal elements, we proceed analogously.
\end{proof}

\begin{theorem}[Dilworth's theorem]\label{thm:dilworths_theorem}\mcite[thm. 2.5.2; thm. 2.5.6]{Harzheim2005OrderedSets}
  For a \hyperref[def:partially_ordered_set]{partially ordered set} \( (P, \leq) \) and a natural number \( m \), the \hyperref[def:partial_order_connectedness/width]{width} of \( P \) is \( m \) if and only if it \( P \) can be partitioned into \( m \) \emph{nonempty} \hyperref[def:partial_order_connectedness/chain]{chains}.
\end{theorem}
\begin{comments}
  \item When partitioning \( P \), we implicitly assume that elements of different chains are incomparable.
  \item The \enquote{sufficiency} part of this theorem is due to \textcite{Dilworth1950DecompositionTheorem}, however our formulation and proof is based on \textcite[\S 2.5]{Harzheim2005OrderedSets}.
  \item This theorem is dual to \fullref{thm:mirsky_theorem}.
\end{comments}
\begin{proof}
  \SufficiencySubProof Suppose that \( P \) has width \( m \), i.e. \( m \) is the maximal cardinality of antichains in \( P \). We will partition \( P \) into \( m \) chains.

  \SubProof*{Proof when \( P \) is finite} Denote by \( n \) the cardinality of \( P \). If \( n \) is zero, then so is \( m \), and there is nothing to prove. Suppose that the inductive hypothesis holds for ordered sets of cardinality \( n \), and suppose that \( P \) has \( n + 1 \) elements.

  Let \( \mscrA \) be the family of antichains of length \( m \). For every such antichain \( A \), define
  \begin{align*}
    I_A &\coloneqq \set{ x \given \qexists* {a \in A} x \leq a } = \bigcup_{a \in A} P_{\leq a}, \\
    F_A &\coloneqq \set{ x \given \qexists* {a \in A} x \geq a } = \bigcap_{a \in A} P_{\geq a}.
  \end{align*}

  Any \( x \) in \( P \) must be comparable with at least one \( a \) in \( A \) because otherwise \( A \cup \set{ x } \) would be a larger antichain. Then \( x \) should belong to \( I_A \) or \( F_A \) or possibly both. In the latter case, we have \( x \geq a \) and \( x \leq a' \) for some \( a \) and \( a' \) from \( A \), so, by transitivity, \( a \leq a' \). Since \( A \) is an antichain, \( a \) and \( a' \) should either be incomparable or coincide, so \( a = x = a' \).

  We conclude that \( P = I_A \cup F_A \) and \( A = I_A \cap F_A \).

  Depending on the antichains in \( \mscrA \), we have two cases:
  \begin{itemize}
    \item Suppose that, there exists an antichain \( A_0 \) from \( \mscrA \) such that both \( I_{A_0} \) and \( F_{A_0} \) have at most \( n \) elements.

    By the inductive hypothesis, both \( I_{A_0} \) and \( F_{A_0} \) can be partitioned into families of \( m \) chains. Denote these families by \( \mscrI_{A_0} \) and \( \mscrF_{A_0} \).

    Since \( A_0 = I_{A_0} \cap F_{A_0} \), each member \( a \) of \( A_0 \) belongs to at least one of the chains in \( \mscrI_{A_0} \) and in \( \mscrF_{A_0} \). Furthermore, \( a \) cannot belong to distinct chains \( I \) and \( I' \) because, by replacing \( a \) with the top elements of \( I \) and \( I' \), we would obtain an antichain with \( m + 1 \) elements, contradicting the maximality of \( m \).

    Denote the chain in \( \mscrI_{A_0} \) corresponding to \( a \) by \( I_a \), and similarly for \( \mscrF_{A_0} \) and \( F_a \). To obtain a partition of \( P \) into \( m \) chains, it remains to connect \( I_a \) to \( F_a \) for each \( a \) in \( A \).

    \item Otherwise, for every antichain \( A \) from \( \mscrA \), either \( I_A \) or \( F_A \) coincides with \( P \).

    In the first case, we have
    \begin{equation*}
      P = I_A = \bigcup_{a \in A} P_{\leq a}.
    \end{equation*}

    Every element \( a \) of \( A \) is maximal in \( P \) because \( x \geq a \) implies that \( x \) belongs to \( P_{\leq a'} \) for some \( a' \) in \( A \), and, since \( a \leq a' \), it can only be the case that \( a = x = a' \).

    Conversely, every maximal element \( y \) belongs to \( A \) because \( y \leq a \) for some \( a \) in \( A \), and thus \( a \leq y \) (and \( y = a \)) due to the maximality of \( y \).

    Thus, if \( P = I_A \), then \( A \) is the set of maximal elements of \( P \). Dually, if \( P = F_A \), then \( A \) is the set of minimal elements of \( P \).

    Generalizing on \( A \), we conclude that \( \mscrA \) contains at most two antichains.

    Now fix a maximal chain \( C \) in \( P \). We already saw that every member \( x \) of \( C \) belongs to \( P_{\leq a_x} \) for some maximal element \( a_x \) and to \( P_{\geq b_x} \) for some minimal element \( b_x \). Since \( C \) is a chain, all its elements are comparable, so \( a_x \) and \( b_x \) do not depend on \( x \). We denote these values by \( a_C \) and \( b_C \).

    The antichains in \( P \setminus C \) are also antichains in \( P \), but they necessarily have less than \( m \) elements because they cannot contain \( a_C \) nor \( b_C \).

    Thus, for any antichain \( A_0 \) from \( \mscrA \), the set \( A_0 \setminus \set{ a_C, b_C } \) is an antichain in \( P \setminus C \) of maximal cardinality \( m - 1 \).

    By the inductive hypothesis on \( P \setminus C \), we can partition it into some chains \( C_1, \ldots, C_{m-1} \). Then \( C_1, \ldots, C_{m-1}, C \) is a partition of \( P \) into nonempty chains.
  \end{itemize}

  \SubProof*{Proof when \( P \) is infinite} Denote the \emph{incomparability} relation in \( P \) by \( \parallel \). Since the relation is irreflexive and symmetric, we can regard \( (P, \parallel) \) as a \hyperref[def:undirected_graph]{simple undirected graph}.

  By assumption, \( P \) has width \( m \), and so every finite subset \( F \) has width at most \( m \). By the finitary case of the theorem, \( F \) can be partitioned into \( m \) chains (some possibly empty). This partition gives us an \( m \)-coloring of the graph \( (F, \parallel) \) because \( a \parallel b \) if and only if they belong to different chains.

  \Fullref{thm:erdos_de_bruijn_theorem} then implies that \( (P, \parallel) \) itself has an \( m \)-coloring. The color classes are exactly the chains.

  Furthermore, since each antichain contains at most one element from each color class, and since there is by assumption an antichain of \( m \) elements, all these chains must be nonempty.

  \NecessitySubProof Conversely, let \( \mscrC \) be a partition of \( P \) into \( m \) nonempty chains.

  Using a choice function \( c \) on \( \mscrC \), possibly provided by the \hyperref[def:zfc/choice]{axiom of choice}, we can define the antichain
  \begin{equation*}
    A \coloneqq \set{ c(C) \given C \in \mscrC }.
  \end{equation*}

  This is indeed an antichain because, if we assume that \( c(C) \leq c(C') \) for some chains \( C \) and \( C' \) from \( \mscrC \), then \( \mscrC \) cannot be a partition unless \( C = C' \).

  Furthermore, there is no antichain \( B \) with more than \( m \) elements because, by \fullref{thm:pigeonhole_principle/finitary}, at least one chain in \( \mscrC \) would have to contain comparable elements of \( B \).
\end{proof}

\begin{theorem}[Mirsky's theorem]\label{thm:mirsky_theorem}\mcite[thm. 2]{Mirsky1971DualOfDilworthsTheorem}
  For a \emph{nonempty} \hyperref[def:partially_ordered_set]{partially ordered set} \( (P, \leq) \) and a natural number \( m \), the \hyperref[def:partial_order_connectedness/height]{height} of \( P \) is \( m \) if and only if it \( P \) can be partitioned into \( m + 1 \) \emph{nonempty} \hyperref[def:partial_order_connectedness/antichain]{antichains}.

  If \( P \) is empty, there is a single empty antichain, however the height is undefined.
\end{theorem}
\begin{comments}
  \item When partitioning \( P \), we implicitly assume that elements of different antichains are incomparable.
  \item The \enquote{sufficiency} part of this theorem is due to \textcite[thm. 2]{Mirsky1971DualOfDilworthsTheorem}.
  \item This theorem is dual to \fullref{thm:dilworths_theorem}.
\end{comments}
\begin{proof}
  \SufficiencySubProof We will use induction on \( m \) to show that, if \( P \) has height \( m \), then it can be partitioned into \( m + 1 \) antichains.

  \begin{itemize}
    \item If \( m = 0 \), then \( P \) has a single chain of length \( 0 \), i.e. \( P \) consists of a single element. This single element vacuously forms an antichain.

    \item Suppose that the theorem holds for \( m \), and let \( P \) have height \( m + 1 \).

    By \cref{thm:def:partial_order_connectedness/maximal_element_antichain}, the set \( M \) of maximal elements of \( P \) form an antichain.

    We must show that \( P \setminus M \) cannot have a chain of length \( m + 1 \). Aiming at a contradiction, suppose that such a chain \( C \) exists. Any element of \( C \) is either maximal in \( P \setminus M \) or is comparable with a maximal element, so \( C \) necessarily contains a maximal element of \( P \setminus M \), say \( a \). But then \( a \) is comparable to some maximal element of \( P \), so we can extend \( C \) to a chain of length \( m + 2 \) in \( P \). Since \( P \) is assumed to have height \( m + 1 \), this is not possible.

    So \( P \setminus M \) has height \( m \), and the inductive hypothesis implies that \( P \setminus M \) can be partitioned into \( m + 1 \) nonempty antichains. With \( M \), we obtain a partition of \( P \) into \( m + 2 \) nonempty antichains.
  \end{itemize}

  \NecessitySubProof Conversely, we will use induction on \( m \) to show that, if \( P \) can be partitioned into \( m + 1 \) nonempty antichains, then its height is at most \( m \).

  \begin{itemize}
    \item If \( m = 0 \), then \( P \) is itself an antichain, so any chain in \( P \) consists of a single element, i.e. has height \( m \).

    \item Suppose that the inductive hypothesis holds for \( m \), and that the family \( \mscrC \) partitions \( P \) into \( m + 2 \) nonempty antichains.

    By \cref{thm:def:partial_order_connectedness/maximal_element_antichain}, the set \( M \) of maximal elements of \( P \) form an antichain. Then \( \mscrC \setminus \set{ M } \) is a partition of \( P \setminus M \) into \( m + 1 \) nonempty antichains.

    By the inductive hypothesis, \( P \setminus M \) has chains of length at most \( m \). An element of \( M \) can increase the length of a chain by \( 1 \), so we conclude that the chains in \( P \) have length at most \( m + 1 \). Hence \( P \) has height \( m \).
  \end{itemize}
\end{proof}

\paragraph{Chain conditions}

\begin{definition}\label{def:stabilizing_sequence}\mimprovised
  Consider either the \hyperref[def:order_function/preserving]{ascending sequence}
  \begin{equation*}
    x_1 \leq x_2 \leq x_3 \leq \cdots
  \end{equation*}
  or the \hyperref[def:order_function/reversing]{descending sequence}
  \begin{equation*}
    x_1 \geq x_2 \geq x_3 \geq \cdots
  \end{equation*}

  If there exists an index \( n \) such that \( x_k = x_n \) whenever \( k > n \), we say that the chain \term[bg=стабилизира (\cite[41]{КоцевСидеров2016КомутативнаАлгебра}), ru=стабилизируется (\cite[52]{Яблонский2003ДискретнаяМатематика}), en=stabilizes (\cite[244]{Aluffi2009Algebra})]{stabilizes} at \( n \).
\end{definition}

\begin{definition}\label{def:chain_condition}
  We say that a \hyperref[def:partially_ordered_set]{partially ordered set} \( (P, \leq) \) satisfies the \term[en=ascending chain condition (\cite[def. 2.37(v)]{DaveyPriestley2002LatticeTheory})]{ascending chain condition} (resp. \term[en=ascending chain condition (\cite[def. 2.37(v)]{DaveyPriestley2002LatticeTheory})]{descending chain condition}) if any of the following equivalent conditions hold:
  \begin{thmenum}
    \thmitem{def:chain_condition/maximal}\mcite[180]{Birkhoff1967LatticeTheory} Every nonempty subset of \( P \) has a \hyperref[def:extremal_points/maximal_and_minimal_element]{maximal} (resp. minimal) element.

    \thmitem{def:chain_condition/stabilization}\mcite[def. 2.37(v)]{DaveyPriestley2002LatticeTheory} Every nonstrict ascending (resp. descending) sequence in \( P \) \hyperref[def:stabilizing_sequence]{stabilizes}.

    \thmitem{def:chain_condition/infinite} \( P \) contains no strictly ascending (resp. descending) sequence.

    \thmitem{def:chain_condition/unbounded} \( P \) contains no \hyperref[def:extremal_points/bounds]{unbounded from above} (resp. from below) \hyperref[def:partial_order_connectedness]{chain}.
  \end{thmenum}
\end{definition}
\begin{defproof}
  We will restrict ourselves to ascending sequences, maximal elements and upper bounds.

  \ImplicationSubProof{def:chain_condition/maximal}{def:chain_condition/stabilization} Suppose that every subset of \( P \) has a maximal element.

  Suppose that there is a nonstrict ascending sequence
  \begin{equation*}
    x_1 \leq x_2 \leq x_3 \leq \cdots.
  \end{equation*}

  Then the set \( \set{ x_k \given k \geq 1 } \) has a maximal element, say \( x_{k_0} \). Since \( x_{k_0} \) is maximal, whenever \( k \geq k_0 \), we have \( x_k \geq x_{k_0} \) and hence \( x_k = x_{k_0} \).

  Thus, the sequence stabilizes.

  \ImplicationSubProof{def:chain_condition/stabilization}{def:chain_condition/infinite} Suppose that every ascending sequence in \( P \) stabilizes.

  Every strictly ascending sequence is also nonstrictly ascending, hence it must stabilize, which contradicts our assumption that it is strictly ascending. The obtained contradiction shows that there are not strictly ascending sequences in \( P \).

  \ImplicationSubProof{def:chain_condition/infinite}{def:chain_condition/unbounded} Suppose that \( P \) contains no infinite strictly ascending sequence.

  Suppose also that \( P \) contains an unbounded from above chain \( C \). Let \( x_1 \) be an element of \( C \). Via \fullref{thm:omega_recursion}, we can define a sequence by letting \( x_{n+1} \) be an element of \( C \) strictly larger than \( x_n \) --- \cref{thm:def:partial_order_connectedness/unbounded} guarantees the existence of such an element. Note that the axiom of choice is needed here.

  We have thus constructed a strictly ascending sequence
  \begin{equation*}
    x_1 < x_2 < x_3 < \cdots
  \end{equation*}

  The existence of such a sequence contradicts our assumption, hence also the existence of \( C \).

  \ImplicationSubProof{def:chain_condition/unbounded}{def:chain_condition/maximal} Suppose that \( P \) contains no unbounded from above chain.

  Let \( A \) be some nonempty subset. Suppose that \( A \) has no maximal element. Then we can construct a strictly ascending sequence
  \begin{equation*}
    x_1 < x_2 < x_3 < \cdots
  \end{equation*}

  The set
  \begin{equation*}
    \set{ x_1, x_2, x_3, \ldots }
  \end{equation*}
  is an unbounded from above chain.

  The obtained contradiction shows that \( A \) has a maximal element.
\end{defproof}

\paragraph{Cofinal sets}

\begin{definition}\label{def:cofinal_set}\mcite[141]{Bourbaki2004TheoryOfSets}
  We say that a subset \( A \) of a \hyperref[def:partially_ordered_set]{preordered set} \( (P, \leq) \) is \term{cofinal} if, for every \( x \) in \( P \), there exists some \( y \) in \( A \) such that \( x \leq y \).
\end{definition}

\begin{example}\label{ex:def:cofinal_set}
  We list several examples of \hyperref[def:cofinal_set]{cofinal} and non-cofinal sets.

  \begin{thmenum}
    \thmitem{ex:def:cofinal_set/finite} In a finite set like \( \set{ 0, 1, 2 } \), the set \( \set{ 2 } \) containing the maximum is cofinal. This is generalized by \cref{thm:def:cofinal_set/top} and \cref{thm:def:cofinal_set/maximal}.

    \thmitem{ex:def:cofinal_set/integers} Consider the set \( \BbbZ \) of integers. Clearly the set \( 2\BbbZ \) of even integers is cofinal. This is generalized by \cref{thm:def:totally_ordered_set/cofinal_iff_unbounded}.

    \thmitem{ex:def:cofinal_set/net_convergence} Cofinal sets are encountered in topology when discussing convergence of nets --- see \fullref{sec:net_convergence}.

    \thmitem{ex:def:cofinal_set/regular_cardinals} \hyperref[def:regular_cardinal]{Regular cardinals} are equal to their own \hyperref[def:cofinality]{cofinality}.
  \end{thmenum}
\end{example}

\begin{proposition}\label{thm:def:cofinal_set}
  \hyperref[def:cofinal_set]{Cofinal sets} have the following basic properties:
  \begin{thmenum}
    \thmitem{thm:def:cofinal_set/top} If a preordered set is bounded from above, a subset is cofinal if and only if it contains a \hyperref[def:extremal_points/top_and_bottom]{top element}.

    \thmitem{thm:def:cofinal_set/maximal} A cofinal subset of a \emph{partially ordered set} contains all \hyperref[def:extremal_points/maximal_and_minimal_element]{maximal elements}.

    \thmitem{thm:def:cofinal_set/acc} In a partially ordered set satisfying the \hyperref[def:chain_condition]{ascending chain condition}, if a set contains all \hyperref[def:extremal_points/maximal_and_minimal_element]{maximal elements}, it is cofinal.

    \thmitem{thm:def:cofinal_set/transitive}\mcite[72]{Harzheim2005OrderedSets} For any preordered set \( P \), if \( A \) is cofinal in \( P \) and \( B \) --- in \( A \), then \( B \) is cofinal in \( P \).

    This is a form of \hyperref[def:binary_relation/transitive]{transitivity}.
  \end{thmenum}
\end{proposition}
\begin{proof}
  \SubProofOf{thm:def:cofinal_set/top} Let \( (P, \leq) \) be a preordered set with top element \( \top \).

  \SufficiencySubProof* Let \( A \) be a cofinal subset of \( P \). Then it must contain some element \( x \) such that \( \top \leq x \). By transitivity, for any element \( y \) of \( P \), we have \( y \leq \top \) and \( \top \leq x \), hence \( y \leq x \).

  Then \( x \) is a top element of \( P \) that belongs to \( A \).

  \NecessitySubProof* Suppose that \( \top \) is in \( A \). Fix an arbitrary element \( x \) of \( P \). Since \( \top \) is an upper bound of \( P \), it follows that \( x \leq \top \). But \( \top \) belongs to \( A \) and \( x \) is arbitrary, thus \( A \) is cofinal.

  \SubProofOf{thm:def:cofinal_set/maximal} Suppose that \( A \) is cofinal in \( (P, \leq) \), where \( \leq \) is a partial order. Let \( m \) be a maximal element of \( P \).

  There must exist some \( x \) in \( A \) such that \( m \leq x \). But \( m \) is maximal, thus \( x = m \). Therefore, \( m \) belongs to \( A \).

  Generalizing on \( m \), we conclude that \( A \) contains all maximal elements of \( P \).

  \SubProofOf{thm:def:cofinal_set/acc} Suppose that \( (P, \leq) \) is partially ordered and its subset \( A \) contains all maximal elements.

  Let \( x \) be an arbitrary element of \( P \). Let \( B \) be the set of all elements strictly larger than \( x \).
  \begin{itemize}
    \item If \( B \) is empty, then \( x \) is itself maximal, and by assumption belongs to \( A \).
    \item Otherwise, since \( P \) satisfies the ascending chain condition, \( B \) has a maximal element \( m \). If \( y \) is any element of \( P \) such that \( m \leq y \), by transitivity \( y \) belongs to \( B \) and, since \( m \) is maximal, \( m = y \). Hence, \( m \) is also maximal for \( P \), and thus belongs to \( A \).

    Therefore, \( A \) has an element \( m \) larger than \( x \).
  \end{itemize}

  Generalizing on \( x \), we conclude that \( A \) is cofinal.

  \SubProofOf{thm:def:cofinal_set/transitive} Let \( (P, \leq) \) be any preordered set.

  Fix some \( p \) in \( P \). Since \( A \) is cofinal in \( P \), there exists some \( a \) in \( A \) such that \( p \leq a \), and similarly some \( b \) in \( B \) such that \( a \leq b \). Transitivity of \( \leq \) implies that \( p \leq b \).

  Hence, \( B \) is transitive in \( P \).
\end{proof}

\paragraph{closure operators}

\begin{definition}\label{def:extensive_function}\mcite[111]{Birkhoff1967LatticeTheory}
  We say that an \hyperref[def:function/endofunction]{endofunction} \( f \) on a \hyperref[def:partially_ordered_set]{partially ordered set} \( (P, \leq) \) is \term[en=extensive (\cite[111]{Birkhoff1967LatticeTheory})]{extensive} if \( x \leq f(x) \) for every \( x \) in \( P \).
\end{definition}
\begin{comments}
  \item \Textcite[def. 1.8.5]{Harzheim2005OrderedSets} prefers the term \enquote{extensional operator}, \textcite[186]{DaveyPriestley2002LatticeTheory} prefer \enquote{increasing map} while \textcite[def. 2.12]{Гуров2013ТеорияРешёток} uses \enquote{\textru{рефлексивный оператор}} (\enquote{reflexive operator}).
\end{comments}

\begin{definition}\label{def:idempotent_function}\mcite[205]{UnivalentFoundationsProgram2013HoTT}
  We say that an \hyperref[def:function/endofunction]{endofunction} \( f \) on an arbitrary \hyperref[def:set]{set} \( A \) is \term{idempotent} if \( f(f(x)) = f(x) \) for every \( x \) in \( A \).
\end{definition}
\begin{comments}
  \item Idempotent functions over \( A \) are \hyperref[def:monoid_idempotent]{idempotent elements} in the \hyperref[def:endomorphism_monoid]{endomorphism monoid} over \( A \).
\end{comments}

\begin{definition}\label{def:abstract_closure_operator}\mcite[def. 7.1]{DaveyPriestley2002LatticeTheory}
  Let \( (P, \leq) \) be a \hyperref[def:partially_ordered_set]{partially ordered set}. We say that the function \( \cl: P \to P \) is a \term[ru=оператор замыкания (\cite[def. 4.12]{Гуров2013ТеорияРешёток})]{closure operator} in \( P \) if it is \hyperref[def:extensive_function]{extensive}, \hyperref[def:idempotent_function]{idempotent} and \hyperref[def:order_function/preserving]{order-preserving}.

  We say that \( x \) is \term[ru=замкнутый (\cite[def. 4.12]{Гуров2013ТеорияРешёток})]{closed} with respect to \( \cl \) if \( x = \cl(x) \).
\end{definition}
\begin{comments}
  \item \Cref{thm:closure_operator_minimality} gives an equivalent condition for an element to be closed while \cref{thm:moore_family_closure_operator} simplifies defining closure operators on power sets.

  \item \Textcite[40]{Harzheim2005OrderedSets} and \textcite[146]{DaveyPriestley2002LatticeTheory} call such an operator simply a \enquote{closure operator}, \textcite[def. 26]{Grätzer2011LatticeTheory} uses \enquote{closure system}, \textcite[111]{Birkhoff1967LatticeTheory} uses \enquote{closure operation} but restricts the definition to lattices of subsets.
\end{comments}

\begin{proposition}\label{thm:closure_operator_minimality}
  For a given \hyperref[def:abstract_closure_operator]{closure operator}, we have \( \cl(x) = c \) if and only if \( c \) is the \hyperref[def:extremal_points/maximum_and_minimum]{least} of all closed elements greater than or equal to \( x \).
\end{proposition}
\begin{proof}
  Fix some element \( x \) and consider the set
  \begin{equation*}
    D \coloneqq \set{ d \geq x \given \cl(d) = d }.
  \end{equation*}

  Clearly \( \cl(x) \) belongs to \( D \) because \( \cl \) is extensive.

  \SufficiencySubProof Consider the closure \( \cl(x) \). Let \( d \geq x \) be an arbitrary closed element. We have \( x \leq d \), and since \( \cl \) preserves order, \( \cl(x) \leq \cl(d) = d \).

  Therefore, \( \cl(x) \) is a lower bound of \( D \) that belongs to \( D \), that is, the least element of \( D \).

  \NecessitySubProof Suppose that \( c \) is the least element of \( D \). Then \( x \leq c \leq \cl(x) \). But
  \begin{equation*}
    \cl(x) \leq \underbrace{\cl(c)}_{c} \leq \underbrace{\cl(\cl(x))}_{\cl(x)},
  \end{equation*}
  hence
  \begin{equation*}
    \cl(x) = c.
  \end{equation*}
\end{proof}

\begin{definition}\label{def:moore_family}\mcite[111]{Birkhoff1967LatticeTheory}
  We say that a family of subsets of an arbitrary \hyperref[def:set]{set} is a \term{Moore family} if it is closed under arbitrary (including empty) intersections.
\end{definition}

\begin{proposition}\label{thm:moore_family_closure_operator}
  Let \( X \) be some set and \( \mscrL \) be a \hyperref[def:moore_family]{Moore family} in \( X \). Then the following function is a \hyperref[def:abstract_closure_operator]{closure operator} on \( X \):
  \begin{equation*}
    \begin{aligned}
      &\cl: \pow(X) \to \mscrL, \\
      &\cl(A) \coloneqq \bigcap \set{ L \in \mscrL \given A \subseteq L }. \\
    \end{aligned}
  \end{equation*}
\end{proposition}
\begin{comments}
  \item This proposition implies that \( \cl(A) \) is the intersection of all closed sets containing \( A \).
  \item \Cref{thm:closure_operator_minimality} implies that \( \cl(A) \) is the smallest closed set containing \( A \).
  \item This allows us to introduce a closure operator on arbitrary families that are closed under intersection --- including \hyperref[def:topological_space]{topological closed sets}, \hyperref[def:affine_hull]{affine hulls}, \hyperref[def:convex_hull]{convex hulls} and \hyperref[def:fol_generated_substructure]{generated first-order substructures} (groups, rings, \( R \)-modules and lattices, among others --- see \cref{ex:def:fol_theory_model_functor}).
\end{comments}
\begin{proof}
  \SubProofOf[def:extensive_function]{extensiveness} The intersection \( \cl(A) \) of \( \set{ L \in \mscrL \given A \subseteq L } \) obviously contains \( A \).

  \SubProofOf[def:idempotent_function]{idempotence} Note that \( \cl(A) \) itself belongs to \( \mscrL \), thus
  \begin{equation*}
    \cl(\cl(A)) = \bigcap \set{ L \in \mscrL \given \cl(A) \subseteq L } = \cl(A).
  \end{equation*}

  \SubProofOf[def:order_function/preserving]{monotonicity} If \( A \subseteq B \), then every set from \( \mscrL \) containing \( B \) also contains \( A \), hence \( \cl(A) \subseteq \cl(B) \).
\end{proof}
